\section{Methodology}
\label{sec:method}

\subsection{Task Formulation}

Given a set of $N$ items $S = \{s_1, s_2, \ldots, s_N\}$, the task is to generate a yes/no question $q$ such that the resulting partition $(S_{\text{yes}}, S_{\text{no}})$ maximizes binary entropy.
Let $p = |S_{\text{yes}}| / N$ denote the fraction of items answered ``yes.''
The binary entropy of the split is:
\begin{equation}
    H(p) = -p \log_2 p - (1-p) \log_2 (1-p)
    \label{eq:entropy}
\end{equation}
This is maximized at $H(0.5) = 1.0$ when $|S_{\text{yes}}| = |S_{\text{no}}| = N/2$, corresponding to a perfect half-split.
A question where all items receive the same answer yields $H(0) = H(1) = 0$.

\subsection{Evaluation Pipeline}

For each item set, our evaluation proceeds in three steps:

\para{Step 1: Generate.} We prompt an LLM to produce a single yes/no question about the items in $S$.
The prompt varies by strategy (described below).

\para{Step 2: Annotate.} We use \gptfull as an oracle to label each item $s_i \in S$ as ``yes'' or ``no'' for the generated question $q$.
The judge operates at temperature 0 for deterministic labeling.

\para{Step 3: Measure.} We compute the binary entropy $H(p)$ of the resulting split using \eqref{eq:entropy}, the perfect split rate (whether $|S_{\text{yes}}| = N/2$), and the mean deviation from half ($|{|S_{\text{yes}}|} - N/2|$).

Using \gptfull as a consistent judge across all conditions ensures that differences in measured entropy reflect differences in question quality, not annotation variance.

\subsection{Prompt Strategies}
\label{sec:prompts}

We evaluate three prompting strategies of increasing specificity:

\para{\basic.} The model receives only the item list and is asked to ``generate a single yes/no question about these items.''
No mention is made of splitting, balancing, or information gain.
This tests the model's default question-asking behavior.

\para{\explicit.} The prompt explicitly instructs the model to ``generate a yes/no question such that exactly $N/2$ items would be answered `yes' and $N/2$ would be answered `no'.''
This tests whether the model can follow a direct splitting instruction.

\para{\chainofthought (Chain-of-Thought).} The prompt asks the model to reason step by step: (1) identify properties of each item, (2) find a property shared by exactly half the items, and (3) formulate a question based on that property.
This tests whether structured reasoning improves split quality.

\subsection{Datasets}
\label{sec:datasets}

We evaluate on 8 datasets spanning a range of set compositions and difficulty levels, totaling 395 item sets (\tabref{tab:datasets}).

\begin{table}[t]
    \centering
    \caption{Datasets used in our evaluation. Difficulty is assessed based on the availability of obvious binary splits within each set. Sources: Bertolazzi et al.\ (\citeyear{bertolazzi2023chatgpt}), Zhang et al.\ (\citeyear{zhang2024probing}), Srivastava et al.\ (\citeyear{srivastava2023bigbench}).}
    \begin{tabular}{@{}llcll@{}}
        \toprule
        \textbf{Dataset} & \textbf{Source} & \textbf{Items/Set} & \textbf{Sets} & \textbf{Difficulty} \\
        \midrule
        Mixed Categories & Custom & 8 & 20 & Easy \\
        \mcrae 8 & Bertolazzi 2023 & 8 & 30 & Easy--Medium \\
        GPT 8 & Bertolazzi 2023 & 8 & 20 & Easy--Medium \\
        \mcrae 16 & Bertolazzi 2023 & 16 & 20 & Medium \\
        \things 8 & Zhang 2024 & 8 & 30 & Medium--Hard \\
        Animals 8 & Custom & 8 & 20 & Hard \\
        Fruits 8 & Custom & 8 & 20 & Hard \\
        \bigbench & Srivastava 2023 & 29 & 15 & Variable \\
        \bottomrule
    \end{tabular}
    \label{tab:datasets}
\end{table}

The datasets vary along two key dimensions.
First, \textbf{category heterogeneity}: Mixed Categories sets contain items from two distinct categories (\eg animals and vehicles), providing an obvious binary split.
\mcrae and GPT sets draw from two taxonomic categories, while Animals and Fruits sets contain items from a single category, requiring the model to find subtler distinguishing properties.
Second, \textbf{set size}: most datasets use 8-item sets, \mcrae 16 uses 16 items, and \bigbench uses 29 items per set.

\subsection{Models}

We evaluate two models:
\begin{itemize}[leftmargin=*,itemsep=0pt,topsep=0pt]
    \item \gptfull: OpenAI's state-of-the-art model (March 2025), used as both generator and judge.
    \item \gptmini: A smaller, more efficient model, used as generator only.
\end{itemize}

Generation uses temperature 0.3; judging uses temperature 0.
All experiments use a fixed random seed of 42, with 10 concurrent API calls per batch.
The total evaluation comprises 1,049 trials across all model--strategy--dataset combinations.
